\chapter{วัฏจักรชีวธรณีเคมีและบทบาทของจุลินทรีย์}
\label{chap:ch02}

% ==========================================================
% แผนบริหารการสอนประจำบทที่ 2
% ==========================================================
\section*{แผนบริหารการสอนประจำบทที่ 2}
\addcontentsline{toc}{section}{แผนบริหารการสอนประจำบทที่ 2}

\noindent\textbf{หัวข้อเนื้อหาประจำบท}
\begin{enumerate}
    \item แนวคิดพื้นฐานของวัฏจักรชีวธรณีเคมีในระบบนิเวศการเกษตร
    \item วัฏจักรไนโตรเจนและเอนไซม์วิทยาของการตรึงไนโตรเจนทางชีวภาพ
    \item กระบวนการไนตริฟิเคชัน ดีไนตริฟิเคชัน และการสูญเสียไนโตรเจนในแปลงเกษตร
    \item วัฏจักรคาร์บอนและการย่อยสลายอินทรียวัตถุในดิน
    \item วัฏจักรฟอสฟอรัสและบทบาทของจุลินทรีย์ละลายฟอสเฟต
\end{enumerate}

\noindent\textbf{วัตถุประสงค์เชิงพฤติกรรม}
\begin{enumerate}
    \item อธิบายขั้นตอนและสมการปฏิกิริยาเคมีในวัฏจักรไนโตรเจน คาร์บอน และฟอสฟอรัสในระบบนิเวศเกษตรได้อย่างถูกต้อง
    \item ระบุชนิดและบทบาทของจุลินทรีย์จำเพาะที่ทำหน้าที่ในแต่ละขั้นตอนของวัฏจักรธาตุอาหารในดินได้อย่างครบถ้วน
    \item วิเคราะห์กลไกการละลายฟอสเฟตด้วยกรดอินทรีย์และเอนไซม์ฟอสฟาเทสของจุลินทรีย์ดินได้อย่างถูกต้อง
    \item คำนวณและประมาณปริมาณธาตุอาหารที่ปลดปล่อยจากการย่อยสลายซากพืชในแปลงปลูกได้อย่างถูกต้อง
\end{enumerate}

\noindent\textbf{กิจกรรมการเรียนการสอน}
\begin{enumerate}
    \item การบรรยายเชิงวิชาการร่วมกับการวิเคราะห์แผนภาพเส้นทางการเคลื่อนย้ายธาตุอาหาร
    \item ปฏิบัติการทดสอบการละลายฟอสเฟตของแบคทีเรียบนอาหาร Pikovskaya Agar
    \item การอภิปรายการจัดการปุ๋ยไนโตรเจนเพื่อลดการปล่อยก๊าซไนตรัสออกไซด์ ($N_2O$)
\end{enumerate}

\noindent\textbf{สื่อการเรียนการสอน}
\begin{enumerate}
    \item เอกสารคำสอนบทที่ 2 และแผนผังเวกเตอร์วัฏจักรชีวธรณีเคมี
    \item ชุดทดสอบการละลายแคลเซียมฟอสเฟตและวัดขนาดวงใส (Clear Zone)
\end{enumerate}

\noindent\textbf{การประเมินผล}
\begin{enumerate}
    \item การประเมินผลจากการทำแผนผังความคิด (Concept Mapping) สรุปวัฏจักรธาตุอาหาร
    \item การทดสอบข้อเขียนและการคำนวณสมดุลมวลสารไนโตรเจน
\end{enumerate}

\newpage

% ==========================================================
% เนื้อหาบทเรียน
% ==========================================================

\section{ภาพรวมของวัฏจักรชีวธรณีเคมีในดิน}
\index{วัฏจักรชีวธรณีเคมี (Biogeochemical Cycles)}

ธาตุอาหารที่จำเป็นต่อการเจริญเติบโตของพืชไม่ได้มีอยู่อย่างไม่จำกัด แต่จะถูกหมุนเวียนเปลี่ยนรูปไปมาระหว่างสิ่งไม่มีชีวิต (บรรยากาศ แหล่งน้ำ หิน และดิน) กับสิ่งมีชีวิต (พืช สัตว์ และจุลินทรีย์) ผ่านกระบวนการทางเคมีและชีวภาพที่เรียกว่า \textbf{วัฏจักรชีวธรณีเคมี (Biogeochemical Cycles)}

จุลินทรีย์ดินทำหน้าที่เป็นตัวเร่งปฏิกิริยาทางชีวภาพ (Biological Catalysts) ที่เปลี่ยนรูปธาตุอาหารที่พืชไม่สามารถดูดซึมได้ ให้อยู่ในรูปของสารประกอบอนินทรีย์ที่ละลายน้ำได้ (Plant-Available Inorganic Forms) ผ่านกระบวนการแปรสภาพเป็นแร่ธาตุ (Mineralization) และตรึงธาตุอาหารกลับเข้าไปในชีวมวลผ่านกระบวนการสะสมธาตุอาหาร (Immobilization)

\section{วัฏจักรไนโตรเจนในระบบเกษตรกรรม}
\index{วัฏจักรไนโตรเจน (Nitrogen Cycle)}

ก๊าซไนโตรเจน ($N_2$) มีสัดส่วนถึง 78\% ของปริมาตรบรรยากาศโลก ทว่าพืชไม่สามารถนำ $N_2$ มาใช้ประโยชน์ได้โดยตรง เนื่องจากมีพันธะสามโคเวเลนต์ ($N \equiv N$) ที่แข็งแรงอย่างยิ่ง การนำไนโตรเจนเข้าสู่ห่วงโซ่อาหารจึงต้องอาศัยจุลินทรีย์ที่มีเอนไซม์จำเพาะ

\subsection{การตรึงไนโตรเจนทางชีวภาพ (Biological Nitrogen Fixation: BNF)}
\index{การตรึงไนโตรเจนทางชีวภาพ (BNF)}
กระบวนการตรึงไนโตรเจนทางชีวภาพเป็นปฏิกิริยารีดักชันที่เปลี่ยนก๊าซ $N_2$ ให้กลายเป็นแอมโมเนีย ($NH_3$) โดยมีตัวเร่งปฏิกิริยาคือเอนไซม์คอมเพล็กซ์ \textbf{ไนโตรจีเนส (Nitrogenase)} ซึ่งประกอบด้วยโปรตีนเหล็ก (Fe protein) และโปรตีนโมลิบดีนัม-เหล็ก (MoFe protein) สมการปฏิกิริยาโดยรวมแสดงไว้ดังสมการ
\begin{equation}
N_2 + 8H^+ + 8e^- + 16\text{ATP} \xrightarrow{\text{Nitrogenase}} 2NH_3 + H_2 + 16\text{ADP} + 16P_i
\end{equation}

เนื่องจากเอนไซม์ไนโตรจีเนสจะสูญเสียสภาพการทำงานอย่างถาวรเมื่อสัมผัสกับก๊าซออกซิเจน จุลินทรีย์ตรึงไนโตรเจนจึงต้องมีกลไกปกป้องเอนไซม์ เช่น ในปมรากถั่วของเชื้อ \textit{Rhizobium} จะมีโปรตีนสีแดงชื่อว่า \textbf{เลกฮีโมโกลบิน (Leghemoglobin)} ทำหน้าที่ดักจับออกซิเจนเพื่อรักษาแรงดันย่อยของออกซิเจนให้อยู่ในระดับต่ำมาก แต่ยังเพียงพอต่อการหายใจระดับเซลล์ของแบคทีเรีย

\subsection{กระบวนการไนตริฟิเคชัน (Nitrification)}
\index{ไนตริฟิเคชัน (Nitrification)}
ไนตริฟิเคชันเป็นกระบวนการออกซิเดชันของแอมโมเนียม ($NH_4^+$) กลายเป็นไนไตร์ต์ ($NO_2^-$) และตามด้วยไนเตรต ($NO_3^-$) ภายใต้สภาวะที่มีออกซิเจนเพียงพอ ดำเนินการโดยแบคทีเรียคีโมอโตโทรฟ 2 ขั้นตอน ดังนี้
\begin{enumerate}
    \item \textbf{ขั้นตอนออกซิไดซ์แอมโมเนีย (Ammonia Oxidation)} โดยแบคทีเรีย \textit{Nitrosomonas} ดังสมการ
    \begin{equation}
    2NH_4^+ + 3O_2 \xrightarrow{\text{Nitrosomonas}} 2NO_2^- + 4H^+ + 2H_2O + \text{พลังงาน}
    \end{equation}
    \item \textbf{ขั้นตอนออกซิไดซ์ไนไตร์ต์ (Nitrite Oxidation)} โดยแบคทีเรีย \textit{Nitrobacter} ดังสมการ
    \begin{equation}
    2NO_2^- + O_2 \xrightarrow{\text{Nitrobacter}} 2NO_3^- + \text{พลังงาน}
    \end{equation}
\end{enumerate}

\subsection{กระบวนการดีไนตริฟิเคชัน (Denitrification)}
\index{ดีไนตริฟิเคชัน (Denitrification)}
เมื่อดินเกิดสภาวะน้ำท่วมขังและขาดออกซิเจน แบคทีเรียกลุ่มไม่ใช้ออกซิเจนแบบเลือกได้ (Facultative Anaerobes) เช่น \textit{Pseudomonas stutzeri} และ \textit{Paracoccus denitrificans} จะใช้ไนเตรตเป็นตัวรับอิเล็กตรอนตัวสุดท้ายในการหายใจ ทำให้ไนเตรตถูกรีดิวซ์กลับเป็นก๊าซไนโตรเจนและก๊าซไนตรัสออกไซด์ ($N_2O$) ดังสมการ
\begin{equation}
NO_3^- \to NO_2^- \to NO \to N_2O \to N_2
\end{equation}
ก๊าซไนตรัสออกไซด์เป็นก๊าซเรือนกระจกที่มีศักยภาพก่อให้เกิดภาวะโลกร้อน (GWP) สูงกว่าคาร์บอนไดออกไซด์ถึง 298 เท่า

\begin{figure}[H]
\centering
\begin{tikzpicture}[scale=0.95, node distance=2.2cm]
    % Nodes
    \node (n2) [rectangle, rounded corners=6pt, draw=agriGreen, fill=agriMint, line width=1.5pt, minimum width=3.5cm, minimum height=1.0cm, align=center, font=\footnotesize\bfseries] {ก๊าซไนโตรเจนในบรรยากาศ\\($N_2$)};
    
    \node (organic) [rectangle, rounded corners=6pt, draw=agriEarth, fill=agriEarth!15, line width=1.2pt, minimum width=3.5cm, minimum height=1.0cm, align=center, font=\footnotesize\bfseries, below left=1.8cm and 0.5cm of n2] {สารอินทรีย์ไนโตรเจนในดิน\\(ซากพืช ซากสัตว์ ฮิวมัส)};

    \node (nh4) [rectangle, rounded corners=6pt, draw=agriBlue, fill=agriBlue!15, line width=1.5pt, minimum width=3.0cm, minimum height=1.0cm, align=center, font=\footnotesize\bfseries, below=2.2cm of organic] {แอมโมเนียม\\($NH_4^+$)};

    \node (no2) [rectangle, rounded corners=6pt, draw=agriGold, fill=yellow!20, line width=1.2pt, minimum width=2.5cm, minimum height=1.0cm, align=center, font=\footnotesize\bfseries, right=2.0cm of nh4] {ไนไตรต์\\($NO_2^-$)};

    \node (no3) [rectangle, rounded corners=6pt, draw=agriLeaf, fill=agriLeaf!20, line width=1.5pt, minimum width=3.0cm, minimum height=1.0cm, align=center, font=\footnotesize\bfseries, right=2.0cm of no2] {ไนเตรต\\($NO_3^-$)};

    % Arrows
    % Fixation: N2 -> NH4 or Plant
    \draw[-{Stealth[scale=1.2]}, line width=1.5pt, color=agriGreen] 
        (n2.west) .. controls (-3.5, 0.5) and (-3.5, -3.5) .. (nh4.west)
        node[midway, left, align=center, font=\scriptsize\bfseries\color{agriGreen}] {การตรึงไนโตรเจน (BNF)\\(\textit{Rhizobium}, \textit{Azotobacter})};

    % Mineralization: Organic -> NH4
    \draw[-{Stealth[scale=1.2]}, line width=1.2pt, color=agriEarth] 
        (organic) -- (nh4)
        node[midway, right, align=left, font=\scriptsize\color{agriEarth}] {แอมโมนิฟิเคชัน\\(Ammonification)};

    % Nitrosomonas: NH4 -> NO2
    \draw[-{Stealth[scale=1.2]}, line width=1.4pt, color=agriBlue] 
        (nh4) -- (no2)
        node[midway, above, align=center, font=\scriptsize\bfseries\color{agriBlue}] {\textit{Nitrosomonas}\\+$O_2$};

    % Nitrobacter: NO2 -> NO3
    \draw[-{Stealth[scale=1.2]}, line width=1.4pt, color=agriLeaf] 
        (no2) -- (no3)
        node[midway, above, align=center, font=\scriptsize\bfseries\color{agriLeaf}] {\textit{Nitrobacter}\\+$O_2$};

    % Denitrification: NO3 -> N2
    \draw[-{Stealth[scale=1.2]}, line width=1.5pt, color=agriRed, dashed] 
        (no3.north) .. controls (5.5, 1.0) and (3.0, 1.5) .. (n2.east)
        node[midway, above right, align=center, font=\scriptsize\bfseries\color{agriRed}] {ดีไนตริฟิเคชัน (Denitrification)\\(\textit{Pseudomonas}, \textit{Paracoccus})\\สูญเสียเป็น $N_2O$ และ $N_2$};

    % Plant uptake
    \draw[-{Stealth[scale=1.2]}, line width=1.2pt, color=agriLeaf] 
        (no3.north west) .. controls (4.0, -1.0) .. (organic.south east)
        node[midway, below, font=\scriptsize\color{agriLeaf}] {พืชดูดซับไปสังเคราะห์โปรตีน};
\end{tikzpicture}
\caption{แผนผังวัฏจักรไนโตรเจนสมบูรณ์ในดินเกษตรกรรม (Complete Soil Nitrogen Cycle)}
\label{fig:nitrogen_cycle}
\end{figure}

\section{วัฏจักรฟอสฟอรัสและจุลินทรีย์ละลายฟอสเฟต}
\index{จุลินทรีย์ละลายฟอสเฟต (Phosphate-Solubilizing Microorganisms)}

ฟอสฟอรัส (P) เป็นธาตุอาหารหลักที่จำเป็นต่อการพัฒนาของระบบรากและการถ่ายโอนพลังงานในรูป ATP อย่างไรก็ตาม ฟอสฟอรัสในดินมากกว่า 95\% มักอยู่ในรูปที่ไม่ละลายน้ำ เช่น เกิดการจับตัวกับเหล็ก ($FePO_4$) หรืออะลูมิเนียม ($AlPO_4$) ในดินกรด และจับกับแคลเซียม ($Ca_3(PO_4)_2$) ในดินด่าง

จุลินทรีย์ละลายฟอสเฟต (Phosphate-Solubilizing Microorganisms) เช่น แบคทีเรีย \textit{Bacillus megaterium}, \textit{Pseudomonas putida} และเชื้อรา \textit{Aspergillus niger} สามารถละลายฟอสเฟตที่ถูกตรึงไว้ได้ผ่าน 2 กลไกสำคัญ
\begin{enumerate}
    \item \textbf{การหลั่งกรดอินทรีย์ (Organic Acid Secretion)} จุลินทรีย์ผลิตกรดกลูโคนิก (Gluconic acid) กรดซิตริก (Citric acid) และกรดออกซาลิก (Oxalic acid) เพื่อลดค่า pH และทำปฏิกิริยาคีเลชัน (Chelation) กับไอออนโลหะ ($Fe^{3+}, Al^{3+}, Ca^{2+}$) ทำให้ปลดปล่อยออร์โธฟอสเฟต ($H_2PO_4^-$ หรือ $HPO_4^{2-}$) ออกสู่สารละลายดิน
    \item \textbf{การหลั่งเอนไซม์ฟอสฟาเตส (Phosphatase Enzymes)} ย่อยสลายสารประกอบฟอสฟอรัสอินทรีย์ เช่น ไฟเตต (Phytate) และกรดนิวคลีอิก
\end{enumerate}

\subsection{กรณีศึกษาวิจัย แบคทีเรียละลายฟอสเฟตทนร้อนในพื้นที่ปกปักทรัพยากรและดินสวนผลไม้}
\index{แบคทีเรียละลายฟอสเฟตทนร้อน}

ในสภาพแปลงเกษตรกรรมเขตร้อนของประเทศไทย อุณหภูมิผิวดินในช่วงฤดูแล้งอาจสูงเกิน 40 องศาเซลเซียส ซึ่งส่งผลให้จุลินทรีย์ดินทั่วไปสูญเสียกิจกรรม การศึกษาของ \textbf{จิรภัทร จันทมาลี และ วาสนา ภักดี (Chanthamalee \& Puckdee, 2017)} เรื่องการคัดแยกแบคทีเรียละลายฟอสเฟตทนร้อนจากพื้นที่ปกปักทรัพยากร มหาวิทยาลัยราชภัฏรำไพพรรณี (อพ.สธ.) และการคัดเลือกแบคทีเรียละลายฟอสเฟตจากดินสวนผลไม้ในจังหวัดจันทบุรี \textbf{(Chanthamalee et al., 2019)} ได้ค้นพบสายพันธุ์แบคทีเรียทนร้อนที่มีศักยภาพสูง

ผลการทดลองพบว่า แบคทีเรียกลุ่มสร้างสปอร์ทนร้อนสกุล \textit{Bacillus} ที่คัดแยกได้สามารถเจริญและสร้างวงใสในการละลายไตรแคลเซียมฟอสเฟต ($Ca_3(PO_4)_2$) บนอาหารเลี้ยงเชื้อ Pikovskaya ที่อุณหภูมิสูงถึง 45 องศาเซลเซียส โดยมีค่าดัชนีการละลายฟอสเฟต (Solubilization Index: SI) สูงกว่า 2.5 ร่วมกับการลดลงของค่า pH ในน้ำเลี้ยงเชื้อจาก 7.0 เหลือต่ำกว่า 4.2 เนื่องจากการหลั่งกรดกลูโคนิกและกรดแลกติก การค้นพบนี้ชี้ให้เห็นศักยภาพอันโดดเด่นในการนำสายพันธุ์แบคทีเรียทนร้อนประจำถิ่นไปต่อยอดเป็นหัวเชื้อปุ๋ยชีวภาพสำหรับสวนผลไม้เมืองร้อน

\begin{agribox}{การใช้ปุ๋ยชีวภาพไมคอร์ไรซาร่วมกับหินฟอสเฟต}
การใช้เชื้อราอาร์บัสคูลาร์ไมคอร์ไรซา (AMF) ร่วมกับหินฟอสเฟตบดละเอียดและจุลินทรีย์ละลายฟอสเฟต เป็นแนวทางเพิ่มประสิทธิภาพการใช้ฟอสฟอรัสในสวนไม้ผลเมืองร้อน เส้นใยราไมคอร์ไรซาจะช่วยขยายพื้นที่ดูดซับฟอสเฟตที่จุลินทรีย์ละลายออกมา ส่งผลให้เกษตรกรสามารถลดการใช้ปุ๋ยเคมีฟอสเฟตสังเคราะห์ลงได้ถึง 30--50\%
\end{agribox}

\section*{คำถามทบทวนท้ายบทที่ 2}
\addcontentsline{toc}{section}{คำถามทบทวนท้ายบทที่ 2}
\begin{enumerate}
    \item จงอธิบายกลไกการทำงานของเอนไซม์ไนโตรจีเนส และเหตุใดเลกฮีโมโกลบินจึงมีความสำคัญอย่างยิ่งในปมรากถั่ว
    \item สภาวะน้ำท่วมขังในแปลงปลูกส่งผลต่อกระบวนการไนตริฟิเคชันและดีไนตริฟิเคชันอย่างไร
    \item จงวิเคราะห์กลไกทางเคมีที่จุลินทรีย์ใช้กรดอินทรีย์ในการละลายฟอสเฟตที่ถูกตรึงในดินกรด
    \item เพราะเหตุใดคุณสมบัติการทนร้อน (Thermotolerance) ของแบคทีเรียละลายฟอสเฟตจึงมีความสำคัญต่อการผลิตปุ๋ยชีวภาพในเขตภูมิอากาศร้อนชื้น
\end{enumerate}

\clearpage